\documentclass{jarticle}
\usepackage[dvipdfmx]{graphicx}
\usepackage{amsmath,amssymb}
\topmargin       -10.0mm
\oddsidemargin   -0.4mm
\evensidemargin  -0.4mm
\headheight       0.0mm
\headsep          0.0mm
\footskip        10.0mm
%
\textheight     260.0mm
\textwidth      170.0mm
%
\setlength{\parindent}{1zw}
%
\usepackage{multirow}
\usepackage{wrapfig}
\usepackage{float}
\makeatletter
\def \fgcaption{\def \@captype{figure} \caption}
\def \tbcaption{\def \@captype{table} \caption}
\makeatother
\def \figurename{Fig.}
\def \tablename{Table}
%
\def\topfraction{.99}
\def\bottomfraction{.99}
\def\textfraction{.01}
\def\floatpagefraction{.99}
%
\pagestyle{empty}
%
\begin{document}
%
\begin{center}
 \begin{tabular}{|c|c|c|} \hline
  \parbox{50mm}{\vspace{1mm} \centering{\Large\bf 2026/11/26}} &
  \parbox{55mm}{\vspace{1mm} \centering{\Large\bf B4　堀田裕樹}} &
  \parbox{50mm}{\vspace{1mm} \centering{\Large\bf No.~16}}
  \\ [2mm] \hline
%
  \multicolumn{3}{|c|}{%
   \parbox{155mm}{\centering\Large\bf
    \vspace{1mm}
    サロゲートモデル構築：精度評価と適応的サンプリング
   }
  }
  \\ [2mm] \hline\hline
%
  \multicolumn{3}{|c|}{%
   \parbox{155mm}{\large\vspace{1mm}
%
    {\Large\bf 今週の進捗要点}
%
    \begin{itemize}
     \item 試験体の範囲だけに配置したテストデータ80点を追加で計算し，逆推定に使う範囲で GPR の精度を評価した
     \item 誤差の大きい領域に教師データ20点を適応的に追加し，200点で学習し直した
     \item 学習・評価に使う周波数帯を変えて予測誤差を比べ，誤差が高周波側に集中することを確かめた
    \end{itemize}
%
   }
  }
  \\ [2mm] \hline
 \end{tabular}
\end{center}
\par



\section{試験体の範囲での精度評価}
試験体の範囲だけに LHS で80点を配置し，FDTD で計算した．

誤差は，溝ありのスペクトル $D$ の相対誤差
\begin{equation}
 \varepsilon_D = \frac{\| D_{\mathrm{GPR}} - D_{\mathrm{FDTD}} \|}{\| D_{\mathrm{FDTD}} \|}
\end{equation}
で評価し，実験で同じ条件を繰り返し測ったときの $D$ のばらつき（中央値 0.078，以下「実験ノイズ」）と比べた．

\section{適応的サンプリング}
誤差の大きい領域に教師データを足すため，候補点ごとに予想誤差
\begin{equation}
 s(\boldsymbol{x}) = \max \left( \sigma_D(\boldsymbol{x}),\ \bar{\varepsilon}(\boldsymbol{x}) \right)
\end{equation}
を求めた．$\sigma_D$ は GPR の予測の標準偏差を $D$ の相対誤差に換算したもの，
$\bar{\varepsilon}$ は近くの学習点5点で，その点を学習から除いて予測したときの誤差を距離で重みづけて平均したものである．
$\sigma_D$ だけでは実際の誤差を小さく見積もる場合があるため，実際の予測誤差から見積もった $\bar{\varepsilon}$ と大きい方をとった．
候補は誤差の大きい角を囲む $p$ 1.22--1.38~mm，$w$ 0.15--0.40~mm，$d$ 0.15--0.30~mm の格子点 7043 点とし，
$s$ の大きい順に，既存の学習点とすでに選んだ点のいずれからも，入力空間（各変数を $[0,1]$ に正規化）で 0.12 以上離れた点だけを採用し，20点を選んだ．
選ぶ前の時点で，試験体の範囲のうち予想誤差が実験ノイズを超える領域は 15.2\% であった．
選んだ20点を FDTD で計算し，教師データを200点にして学習し直した．

\section{学習し直した結果}
180点と200点の GPR の誤差を Table~\ref{tab:retrain} にまとめる．
実験の試験体12条件では，$D$ ではなく $R = D / D_{\mathrm{smooth}}$ の相対誤差 $\varepsilon_R$ で比べ，
実験の $R$ のばらつき（平滑面の測定のばらつき込み，中央値 0.120）を基準とした．

\begin{table}[H]
 \centering
 \caption{180点と200点の GPR の誤差}
 \label{tab:retrain}
 \begin{tabular}{llcc} \hline
  評価に使う点 & 指標 & 180点 & 200点 \\ \hline
  試験体の範囲のテスト80点 & $\varepsilon_D$ の中央値／最大 & 0.034／0.099 & 0.029／0.093 \\
                         & 実験ノイズ 0.078 未満 & 77／80 & 79／80 \\
  実験の試験体12条件       & GPR と FDTD の $\varepsilon_R$ の最大 & 0.139 & 0.097 \\
                         & 実験の $R$ のばらつき 0.120 未満 & 11／12 & 12／12 \\ \hline
 \end{tabular}
\end{table}

\begin{figure}[H]
 \centering
 \includegraphics[width=16cm]{../../vscode/surrogate/tmp_output/gpr/train_ph14_adapt1_refpeak_J17_invp_duty/specimen_range_test_spec/error_vs_param_test.png}
 \caption{200点で学習した GPR の試験体の範囲のテスト80点での $\varepsilon_D$ と $p$, $w$, $d$ の関係（破線は実験ノイズ 0.078）}
 \label{fig:err200}
\end{figure}

\subsection*{課題}
\begin{itemize}
 \item 弱点だった角の2条件（$p=1.25$, $w=0.40$, $d=0.15$ および $0.20$~mm）で，
       GPR と FDTD の差が予測の $\pm 2\sigma$ に入る周波数点の割合（理想 0.95）が，0.72 から 0.60，0.76 から 0.62 に下がった．
       誤差は小さくなったが，予測の $\sigma$ がそれ以上に小さくなったためで，GPR が不確かさを小さく見積もりすぎている．
       逆推定で $\sigma$ を誤差として使う前に直す必要がある．
 \item GPR と実験の差が各条件の実験のばらつきを下回ったのは，180点・200点とも12条件中8条件で変わらなかった．
       FDTD と実験の差でも8条件であり，この差は GPR ではなく FDTD と実験の差が支配している．
\end{itemize}


\section{周波数帯による予測誤差の違い}
誤差がどの周波数に集まっているかを調べるため，学習と評価に使う周波数帯を変えて比べた．
本節では教師データ200点で学習したモデルについてのみ論じる．
帯域は平滑面のスペクトルの大きさ（ピーク 4.80~MHz）から次の5つとした．
\begin{itemize}
 \item full：2.00--7.75~MHz（これまでと同じ）
 \item $-6$~dB：3.00--6.60~MHz，\ $-10$~dB：2.55--7.05~MHz
 \item low：2.00--5.30~MHz，\ high：4.30--7.75~MHz（ピークで切らず，互いに 0.5~MHz ずつ重ねた）
\end{itemize}
帯域ごとに PCA と GPR を学習し直し，主成分数は full の $J=17$ と同じ累積寄与率（99.457\%）に届く最小の数とした
（$-6$~dB で12，$-10$~dB で15，low で12，high で13）．
評価には $R$ の二乗平均平方根誤差
\begin{equation}
 \mathrm{RMSE} = \sqrt{ \frac{1}{N_f} \sum_{f} \left| R_{\mathrm{GPR}}(f) - R_{\mathrm{FDTD}}(f) \right|^2 }
\end{equation}
を使った（$N_f$ は帯域内の周波数点の数）．前節までの相対誤差とは定義が異なる．

% \begin{figure}[H]
%  \centering
%  \includegraphics[width=16cm]{../../vscode/surrogate/tmp_output/gpr/band_rmse/rmse_vs_params_full_200.png}
%  \caption{full で学習・評価した RMSE と $p$, $d$, $w$ の関係（200点，試験体の範囲のテスト80点）}
%  \label{fig:rmsefull}
% \end{figure}

\begin{figure}[H]
 \centering
 \includegraphics[width=10cm]{../../vscode/surrogate/tmp_output/gpr/band_rmse/band_matrix_200.png}
 \caption{学習帯域と評価帯域の組み合わせごとの RMSE の中央値（200点，試験体の範囲のテスト80点．空欄は評価帯域が学習帯域の外）}
 \label{fig:bandmat}
\end{figure}


\section{全域で学習し低域で評価した GPR の誤差}
3節の Fig.~\ref{fig:err200} と同じ形で，全域（full，2.00--7.75~MHz）で学習した200点の GPR を低域（low，2.00--5.30~MHz）だけで評価した．
誤差は式 (1) の和を低域の周波数点だけでとった
\begin{equation}
 \varepsilon_D^{\mathrm{low}} = \frac{\| D_{\mathrm{GPR}} - D_{\mathrm{FDTD}} \|_{\mathrm{low}}}{\| D_{\mathrm{FDTD}} \|_{\mathrm{low}}}
\end{equation}
とし，実験ノイズも同じ低域で計算し直した（中央値 0.067）．全域の 0.078 をそのまま使うと，帯域によって基準が変わるためである．

\begin{figure}[H]
 \centering
 \includegraphics[width=16cm]{../../vscode/surrogate/tmp_output/gpr/train_ph14_adapt1_refpeak_J17_invp_duty/specimen_range_test_spec/error_vs_param_eval_low.png}
 \caption{全域で学習し低域で評価した GPR の $\varepsilon_D^{\mathrm{low}}$ と $p$, $w$, $d$ の関係（200点，試験体の範囲のテスト80点．破線は低域の実験ノイズ 0.067）}
 \label{fig:errlow}
\end{figure}

\begin{table}[H]
 \centering
 \caption{評価帯域による GPR の誤差の違い（全域で学習した200点，試験体の範囲のテスト80点）}
 \label{tab:errlow}
 \begin{tabular}{lcc} \hline
  評価帯域 & 全域（2.00--7.75~MHz） & 低域（2.00--5.30~MHz） \\ \hline
  $\varepsilon_D$ の中央値／最大 & 0.029／0.093 & 0.023／0.102 \\
  実験ノイズ（同じ帯域）の中央値 & 0.078 & 0.067 \\
  実験ノイズ未満 & 79／80 & 79／80 \\ \hline
 \end{tabular}
\end{table}



\end{document}
